\scanpage{58}

Rösrath, le 29 octobre 1966\\

\margin{\unsure{original}}

Cher Grothendieck,\\

Je te signale tout d'abord mon adresse en Allemagne, maintenant définitive (jusqu'au $1^{\underline{\mathrm{er}}}$ juillet 67)

S.M.\quad Pierre Deligne\\
Athénée royal\\
BPS 8\\
FBA (forces belges en Allemagne)\\

Je suis en train de lire EGA IV 19 que tu m'as envoyé.

\uline{a) remarque sur les hypothèses de séparation utilisées}

L'hypothèse $(*)$ ``tout quotient de tout sous-module de $M$ \deleted{\unsure{est}} (\deleted{resp} $M^n$) est séparé pour la topologie $J$-adique'' me semble idiote :

(i) elle \deleted{n'implique pas} n'est pas impliquée par l'hypothèse de finitude relative noethérienne standart : ``$M$ fini sur la $A$-algèbre noethérienne $B$, $A$ noethérien, et $JB\subset R(B)$'' qui elle est naturelle.

\margin{?\ \rule{0.25pt}{1.4cm}}

\margin{\uline{Err}}

(ii) le seul exemple non trivial indiqué où cette hypothèse est vérifiée, soit 19.5.6 (ii), est canulé, car \deleted{$(*)$} dans $(*)$ on ne se limite pas à prendre des quotients par des sous-modules gradués.

% slanted
\margin{\uline{\unsure{Ça va pas|ça va p.|ja \illegible{1}}}\ $\Big($}

\margin{?}

[exemple : $A=B[T]$, $M=N[T]$, $\deleted{\unsure{M}}N$ est un $A$-module, $T$ agissant par $1$, et ce $A$-module est quotient de $M$ ]

Je crois qu'il faut éviter d'employer des hypothèses de séparation qu'on ne rencontre pas dans la nature, comme ``idéalement séparé'' etc., et se limiter à utiliser (i) - J'avoue toutefois, avec honte, que dans le même geste sacrilège je serais prêt à supposer tous les schémas séparés et noethériens...

Sinon, pour coiffer (\unsure{i|1}) et $(*)$ resp. on pourrait prendre :

\margin{!\ $\Bigg($}

\[
\begin{aligned}
&\forall f:\ N\longrightarrow \deleted{\unsure{M}}P,\quad \text{si } N \text{ et } \deleted{\unsure{M}}P \text{ sont finis sur } A,\quad \operatorname{Ker}(M\otimes N\deleted{)}\longrightarrow\deleted{\operatorname{Ker}}\deleted{(}M\otimes P)\quad\text{est}\\
&\quad J\text{-séparé.}
\end{aligned}
\]
% archive sheet number: 57
